\section{Feynman 图}

\subsection{微扰展开的图示}

把相互作用 $H_{\mathrm{int}}$ 按 Dyson 级数展开后，每一项用 Wick 定理变成 $G_0$ 线与相互作用顶点的乘积。用图表示：
\begin{itemize}
  \item 有向实线：自由传播子 $G_0$；
  \item 虚线或波浪线：相互作用（瞬时库仑线、声子线等）；
  \item 顶点：守恒量（电荷、动量、自旋）流入流出平衡。
\end{itemize}
每个图对应一个可计算的积分/求和表达式，并带有组合因子与费米子符号。

\subsection{连通图与真空图}

热力学势与基态能移主要由连通真空图贡献；单粒子格林函数的修正来自带有两个外腿的连通图。忽略非连通部分可避免对体积因子的重复计数，并使“强度量”的图求和结构更清晰。

\subsection{连锁集群定理}

有限温度微扰展开中，配分函数
\begin{equation}
  Z=Z_0
  \Bigl\langle
  T_\tau\exp\Bigl(-\int_0^\beta H_{\mathrm{int}}(\tau)\,\mathrm{d}\tau\Bigr)
  \Bigr\rangle_0
\end{equation}
的图求和包含不连通真空泡。\textbf{连锁集群定理}（linked-cluster theorem）断言
\begin{equation}
  \ln(Z/Z_0)
  =\text{全部连通真空图之和},
\end{equation}
从而巨势 $\Omega=-k_BT\ln Z$ 对体积严格可加。这是施均仁讲义微扰论章节的核心结构定理之一。

\subsection{Hugenholtz 图}

对密度--密度型瞬时相互作用，常把直接与交换顶点合并，用粗点表示反对称化相互作用，得到 Hugenholtz 图：图更少、对称因子更清晰，与标准 Feynman 图一一对应。

\subsection{自能图}

单粒子不可约图的和定义为自能 $\Sigma$。最低阶（一阶）对密度-密度相互作用给出 Hartree 与 Fock 贡献：
\begin{itemize}
  \item Hartree：闭合密度泡与另一条传播线通过相互作用相连；
  \item Fock：交换型连接，体现 Pauli 原理。
\end{itemize}
更高阶包含屏蔽、顶点修正等过程。

\begin{example}[二阶自能示意]
对短程相互作用，二阶自能常含“日出”（sunset）图等，积分形式示意性地写作
\begin{equation}
  \Sigma^{(2)}(k)
  \sim
  \sum_{k_1k_2}
  |V|^2\,
  G_0(k_1)G_0(k_2)G_0(k-k_1+k_2),
\end{equation}
具体动量/频率标记依所用绘景（实频或 Matsubara）而定。
\end{example}

\subsection{读图的物理直觉}

图不仅是计算工具，也是过程叙事：哪条线表示粒子、哪条表示空穴、何时发射玻色子、何时发生交换散射。训练读图能力，有助于判断哪些过程在低能或长波极限下主导。
